\chapter{1977年江苏扬州地区初试卷}

\section{解答题}

\begin{problem}
计算：\(8^{\frac{2}{3}}+\left(\frac{1}{2}\right)^{-2}-\log_{2}1+\log_{3}3-\left(\sqrt{2}-1\right)^{0}\).
\end{problem}

\begin{solution}
因为 \(8^{\frac{2}{3}}=\left(\sqrt[3]{8}\right)^2=4\)，\(\left(\frac{1}{2}\right)^{-2}=2^2=4\)，\(\log_{2}1=0\)，\(\log_{3}3=1\)，且 \(\sqrt{2}-1\neq0\)，所以 \(\left(\sqrt{2}-1\right)^0=1\). 因此原式 \(=4+4-0+1-1=8\).
\end{solution}

\begin{problem}
求直线 \(\sqrt{3}x+y-2=0\) 的斜率，倾斜角及在 \(x\) 轴，\(y\) 轴上的截距.
\end{problem}

\begin{solution}
将直线方程化为斜截式，得 \(y=-\sqrt{3}x+2\)，所以直线的斜率为 \(k=-\sqrt{3}\). 设直线的倾斜角为 \(\alpha\)，其中 \(0\leqslant\alpha<\pi\)，则 \(\tan\alpha=k=-\sqrt{3}\)，故 \(\alpha=\frac{2\pi}{3}=120^\circ\).

令 \(y=0\)，得 \(\sqrt{3}x-2=0\)，所以直线与 \(x\) 轴的交点为 \(\left(\frac{2}{\sqrt{3}},0\right)=\left(\frac{2\sqrt{3}}{3},0\right)\)，在 \(x\) 轴上的截距为 \(\frac{2\sqrt{3}}{3}\). 令 \(x=0\)，得 \(y=2\)，所以直线与 \(y\) 轴的交点为 \(\left(0,2\right)\)，在 \(y\) 轴上的截距为 \(2\).
\end{solution}

\begin{problem}
把 \(A\) 点的极坐标 \(\left(4,\frac{3\pi}{4}\right)\) 化成直角坐标；把 \(B\) 点的直角坐标 \(\left(-\sqrt{3},1\right)\) 化成极坐标.
\end{problem}

\begin{solution}
对点 \(A\)，由极坐标与直角坐标的关系 \(x=r\cos\theta\)，\(y=r\sin\theta\)，得 \(x=4\cos\frac{3\pi}{4}=-2\sqrt{2}\)，\(y=4\sin\frac{3\pi}{4}=2\sqrt{2}\). 因此点 \(A\) 的直角坐标为 \(\left(-2\sqrt{2},2\sqrt{2}\right)\).

对点 \(B\)，其极径为 \(r=\sqrt{x^2+y^2}=\sqrt{\left(-\sqrt{3}\right)^2+1^2}=2\). 又因为 \(\cos\theta=\frac{x}{r}=-\frac{\sqrt{3}}{2}\)，\(\sin\theta=\frac{y}{r}=\frac{1}{2}\)，所以 \(\theta=\frac{5\pi}{6}+2k\pi\)，其中 \(k\in\Z\). 取 \(0\leqslant\theta<2\pi\) 时，\(\theta=\frac{5\pi}{6}\)，故点 \(B\) 的一个极坐标为 \(\left(2,\frac{5\pi}{6}\right)\).
\end{solution}

\begin{problem}
不画图象求出函数 \(y=2x^2-4x-1\) 的图象的顶点坐标，对称轴方程，并说明 \(x\) 取何值时，函数取得极大值或极小值？这个极值是多少？
\end{problem}

\begin{solution}
配方得 \(y=2\left(x^2-2x\right)-1=2\left[\left(x-1\right)^2-1\right]-1=2\left(x-1\right)^2-3\). 因为 \(\left(x-1\right)^2\geqslant0\)，所以函数图象的顶点为 \(\left(1,-3\right)\)，对称轴方程为 \(x=1\). 当且仅当 \(x=1\) 时，\(\left(x-1\right)^2=0\)，函数取得极小值 \(-3\). 由于二次项系数 \(2>0\)，函数值向上无界，所以函数没有极大值.
\end{solution}

\begin{problem}
已知 \(x\) 是正的纯小数，化简
\(\sqrt{x^2-2+\frac{1}{x^2}} \cdot \left(x+\frac{1}{x}\right) \div \frac{2x^{-2}}{1-\frac{x^2}{1+x^2}}\).
\end{problem}

\begin{solution}
因为 \(x\) 是正的纯小数，所以 \(0<x<1\)，从而 \(\frac{1}{x}-x>0\). 因此
\[
\sqrt{x^2-2+\frac{1}{x^2}}=\sqrt{\left(\frac{1}{x}-x\right)^2}=\frac{1}{x}-x
\]
又
\[
1-\frac{x^2}{1+x^2}=\frac{1}{1+x^2}
\]
所以
\[
\frac{2x^{-2}}{1-\frac{x^2}{1+x^2}}=\frac{2x^{-2}}{\frac{1}{1+x^2}}=\frac{2\left(1+x^2\right)}{x^2}
\]
于是原式
\[
=\left(\frac{1}{x}-x\right)\left(x+\frac{1}{x}\right)\div\frac{2\left(1+x^2\right)}{x^2}
=\frac{1-x^2}{x}\cdot\frac{1+x^2}{x}\cdot\frac{x^2}{2\left(1+x^2\right)}
=\frac{1-x^2}{2}
\]
\end{solution}

\begin{problem}
在 \(\triangle ABC\) 中，\(\angle A=60^\circ\)，\(\angle B=45^\circ\)，\(BC\) 边上的高 \(AD=10\)，求 \(\triangle ABC\) 的面积（不查表，得数可保留根式）.
\end{problem}

\begin{solution}
因为 \(AD\perp BC\)，所以 \(\triangle ABD\) 和 \(\triangle ACD\) 都是直角三角形. 在直角三角形 \(ABD\) 中，\(\angle ABD=\angle B=45^\circ\)，故 \(\angle BAD=45^\circ\)，从而 \(BD=AD=10\).

又因为 \(\angle C=180^\circ-60^\circ-45^\circ=75^\circ\)，在直角三角形 \(ACD\) 中，\(\tan75^\circ=\frac{AD}{CD}\). 由 \(30^\circ\) 角的直角三角形可知 \(\tan30^\circ=\frac{1}{\sqrt{3}}\)，所以不查表而由正切加法公式得
\[
\tan75^\circ=\tan\left(45^\circ+30^\circ\right)=\frac{1+\frac{1}{\sqrt{3}}}{1-\frac{1}{\sqrt{3}}}=2+\sqrt{3}
\]
因此 \(CD=\frac{10}{2+\sqrt{3}}=10\left(2-\sqrt{3}\right)\)，从而 \(BC=BD+DC=10\left(3-\sqrt{3}\right)\). 所以
\[
S_{\triangle ABC}=\frac{1}{2}BC\cdot AD=\frac{1}{2}\cdot10\left(3-\sqrt{3}\right)\cdot10=50\left(3-\sqrt{3}\right)
\]
\end{solution}

\begin{problem}
圆台上、下底面半径分别是 \(2 \mathrm{~cm}\) 和 \(6 \mathrm{~cm}\)，它的侧面积是 \(40\pi \mathrm{~cm}^2\)，求圆台体积.（\(V\text{体}=\frac{1}{3}\pi h\left(R^2+r^2+Rr\right)\)）
\end{problem}

\begin{solution}
设圆台的母线长为 \(l\)，上、下底面半径分别为 \(r=2 \mathrm{~cm}\)，\(R=6 \mathrm{~cm}\). 由圆台侧面积公式 \(S_{\text{侧}}=\pi\left(R+r\right)l\)，得
\[
40\pi=\pi\left(6+2\right)l
\]
所以 \(l=5 \mathrm{~cm}\). 在圆台的轴截面内，从上底的一端向下底作垂线，得到斜边为 \(l\)、两直角边分别为圆台高 \(h\) 和两底半径之差 \(R-r\) 的直角三角形，因此
\[
h=\sqrt{l^2-\left(R-r\right)^2}=\sqrt{5^2-\left(6-2\right)^2}=3 \mathrm{~cm}
\]
代入题给体积公式，得
\[
V=\frac{1}{3}\pi h\left(R^2+r^2+Rr\right)=\frac{1}{3}\pi\cdot3\left(6^2+2^2+6\cdot2\right)=52\pi \mathrm{~cm}^3
\]
\end{solution}

\begin{problem}
某工厂 \(1977\) 年一月份的产值是 \(10 \text{ 万元}\)，如果以后平均每月产值比前一个月增长 \(10\%\)，问到几月份就能完成今年总产值 \(135 \text{ 万元}\) 的任务？（\(\lg 1.1=0.0414\)，\(\lg 2.35=0.3711\)）
\end{problem}

\begin{solution}
设第 \(n\) 个月的产值为 \(a_n\)，则 \(a_1=10\)，公比为 \(1.1\)，所以 \(a_n=10\cdot1.1^{n-1}\). 前 \(n\) 个月的总产值为
\[
S_n=10\frac{1.1^n-1}{1.1-1}=100\left(1.1^n-1\right)
\]
要完成 \(135\text{ 万元}\) 的任务，需满足 \(S_n\geqslant135\)，即 \(1.1^n\geqslant2.35\). 由于 \(1.1>1\)，取常用对数得
\[
n\geqslant\frac{\lg2.35}{\lg1.1}=\frac{0.3711}{0.0414}\approx8.96
\]
因此最小整数 \(n=9\). 一月份是第一个月，第九个月为九月份，所以到九月份可以完成今年总产值的任务.
\end{solution}

\begin{problem}
已知一元二次方程 \(2x^2+4\left(2e-1\right)x+4e^2-1=0\) 有两个相等的实数根，以 \(e\) 作为椭圆的离心率，这个椭圆的中心在原点，焦点在 \(x\) 轴上，半长轴的长是 \(10\)，求这个椭圆的方程.
\end{problem}

\begin{solution}
原一元二次方程有两个相等的实数根，所以判别式为零，即
\[
\Delta=\left[4\left(2e-1\right)\right]^2-4\cdot2\left(4e^2-1\right)
=16\left(2e-1\right)^2-8\left(4e^2-1\right)
=8\left(4e^2-8e+3\right)=0
\]
因为 \(4e^2-8e+3=(2e-1)(2e-3)\)，所以 \(e=\frac{1}{2}\) 或 \(e=\frac{3}{2}\). 椭圆的离心率满足 \(0<e<1\)，故舍去 \(e=\frac{3}{2}\)，取 \(e=\frac{1}{2}\).

设椭圆的半长轴为 \(a=10\)，焦距的一半为 \(c\)，则离心率 \(e=\frac{c}{a}\)，从而 \(c=ae=5\). 由椭圆关系式 \(b^2=a^2-c^2\)，得 \(b^2=10^2-5^2=75\). 因为焦点在 \(x\) 轴上且中心在原点，椭圆方程为
\[
\frac{x^2}{a^2}+\frac{y^2}{b^2}=1
\]
即 \(\frac{x^2}{100}+\frac{y^2}{75}=1\).
\end{solution}

\begin{problem}
已知 \(\triangle ABC\) 的三边 \(a\)，\(b\)，\(c\) 成等差数列，\(a>b>c\)，\(b\) 边的对角是 \(\angle B\). 求证：\(\angle B\) 的度数小于 \(60^\circ\).
\end{problem}

\begin{solution}
因为 \(a\)，\(b\)，\(c\) 成等差数列且 \(a>b>c\)，设公差为 \(d>0\)，则 \(a=b+d\)，\(c=b-d\). 由余弦定理，
\[
\cos B=\frac{a^2+c^2-b^2}{2ac}
\]
而
\[
a^2+c^2-b^2-ac=(b+d)^2+(b-d)^2-b^2-(b+d)(b-d)=3d^2>0
\]
所以 \(a^2+c^2-b^2>ac\). 因为 \(a\)，\(c>0\)，由余弦定理可得 \(\cos B>\frac{1}{2}=\cos60^\circ\). 三角形内角满足 \(0^\circ<B<180^\circ\)，而余弦函数在此区间上递减，故 \(B<60^\circ\).
\end{solution}
